% ============================================================
%  CHAPTER 5 — REINFORCEMENT LEARNING & FEDERATED LEARNING
% ============================================================
\chapter{Reinforcement \& Federated Learning}

\section{Tabular Q-Learning Load Balancer}

\subsection{Motivation}

IoT traffic can change rapidly due to bursts and attack-like behavior. A static buffering or throttling policy is therefore difficult to tune across operating regimes. Reinforcement learning provides an adaptive control mechanism that learns which action is preferable for a given congestion state.

\subsection{State--Action Design}

The controller uses a discrete state space based on queue depth and a small action space of delay multipliers (Table~\ref{tab:rl_space}).

\begin{table}[H]
\centering
\caption{Q-Learning State--Action Space}
\label{tab:rl_space}
\begin{tabular}{@{}lll@{}}
\toprule
\textbf{Component} & \textbf{Values} & \textbf{Description} \\
\midrule
\multicolumn{3}{@{}l}{\textit{State Space}} \\
\quad Low & Queue $< 10$ & Minimal congestion \\
\quad Medium & $10 \leq$ Queue $< 50$ & Moderate load \\
\quad High & $50 \leq$ Queue $< 100$ & Heavy congestion \\
\quad Critical & Queue $\geq 100$ & Imminent overflow \\
\midrule
\multicolumn{3}{@{}l}{\textit{Action Space}} \\
\quad Accelerate & $0.5\times$ & Speed up processing \\
\quad Maintain & $1.0\times$ & No change \\
\quad Throttle & $2.0\times$ & Slow down processing \\
\bottomrule
\end{tabular}
\end{table}

\subsection{Reward and Update Rule}

The immediate reward is defined to penalize large flow durations:
\begin{equation}
r(s, a) = -\text{flowDuration}(s, a)
\end{equation}

Q-values are updated using the standard temporal-difference rule \cite{qlearning}:
\begin{equation}
Q(s, a) \leftarrow Q(s, a) + \alpha \left[r + \gamma \max_{a'} Q(s', a') - Q(s, a)\right]
\end{equation}

\subsection{Integration}

During packet processing, the simulation estimates a congestion state, selects an action using an $\epsilon$-greedy policy, applies the multiplier, and updates the Q-table. Figure~\ref{fig:rl_loop} illustrates this closed feedback loop.

\begin{figure}[H]
\centering
\begin{tikzpicture}[
    block/.style={rectangle, rounded corners=6pt, draw=violet!60, fill=violet!8,
        minimum width=3.2cm, minimum height=1.2cm,
        font=\small\bfseries, align=center, text=violet!80!black},
    arr/.style={-{Stealth[length=3mm]}, thick, black!60},
]
\node[block] (ENV) {Simulation\\Environment};
\node[block, right=4.0cm of ENV] (AGT) {QTableBalancer\\Agent};

\draw[arr] (ENV.north east) -- ++(0,0.8) -| node[above, pos=0.25, font=\scriptsize]{state $s$ (queue depth)} (AGT.north west);
\draw[arr] (AGT.south west) -- ++(0,-0.8) -| node[below, pos=0.25, font=\scriptsize]{action $a$ (delay multiplier)} (ENV.south east);
\draw[arr, dashed, red!50!black] (ENV.east) -- node[above, font=\scriptsize, text=red!60!black]{reward $r = -$flowDuration} (AGT.west);
\end{tikzpicture}
\caption{Closed-loop reinforcement learning: the simulation provides state and reward signals to the QTableBalancer, which selects a delay multiplier action applied to packet processing.}
\label{fig:rl_loop}
\end{figure}

\section{Federated Learning Pipeline}

\subsection{Motivation}

Centralizing raw telemetry can be undesirable due to privacy and bandwidth constraints. Federated Learning mitigates this by training locally at the edge and sharing only model updates.

\subsection{FedAvg Summary}

The platform prototypes Federated Averaging (FedAvg) \cite{fedavg} in a simplified setting:
\begin{enumerate}
    \item the cloud initializes and broadcasts a global model,
    \item each edge node trains locally on its private data partition,
    \item nodes upload model updates (not raw records), and
    \item the cloud aggregates updates to form the next global model.
\end{enumerate}

Figure~\ref{fig:fedavg_round} visualizes a single communication round.

\begin{figure}[H]
\centering
\begin{tikzpicture}[
    cloud/.style={rectangle, rounded corners=6pt, draw=orange!60, fill=orange!8,
        minimum width=4.5cm, minimum height=1.4cm,
        font=\small\bfseries, align=center, text=orange!80!black},
    edge/.style={rectangle, rounded corners=5pt, draw=blue!50, fill=blue!6,
        minimum width=2.8cm, minimum height=1.0cm,
        font=\scriptsize\bfseries, align=center, text=blue!80!black},
    arr/.style={-{Stealth[length=2.8mm]}, thick, black!60},
    feed/.style={-{Stealth[length=2.8mm]}, thick, dashed, red!50!black},
]
% Cloud aggregator
\node[cloud] (CLD) {Cloud Aggregator\\{\tiny FedAvg: weighted average}};

% Edge nodes
\node[edge, below left=2.5cm and 4cm of CLD] (E0) {Edge Node 0\\{\tiny Local SGD}};
\node[edge, below=2.5cm of CLD] (E1) {Edge Node 1\\{\tiny Local SGD}};
\node[edge, below right=2.5cm and 4cm of CLD] (E2) {Edge Node 2\\{\tiny Local SGD}};

% Upload arrows
\draw[arr] (E0.north) -- node[left, font=\tiny, pos=0.4]{$w_0$} (CLD.south west);
\draw[arr] (E1.north) -- node[right, font=\tiny, pos=0.4]{$w_1$} (CLD.south);
\draw[arr] (E2.north) -- node[right, font=\tiny, pos=0.4]{$w_2$} (CLD.south east);

% Broadcast arrows (dashed red, going back down)
\draw[feed] (CLD.south west) ++(-.3,0) -- ++(-0.8,-1.0) node[left, font=\tiny, text=red!60!black, pos=0.5]{$w_{\text{global}}$} |- (E0.east);
\draw[feed] (CLD.south east) ++(.3,0) -- ++(0.8,-1.0) node[right, font=\tiny, text=red!60!black, pos=0.5]{$w_{\text{global}}$} |- (E2.west);

% Round label
\node[above=0.3cm of CLD, font=\scriptsize\itshape, text=black!50] {One communication round};
\end{tikzpicture}
\caption{FedAvg communication round: edge nodes train locally on private data partitions, upload model weights to the cloud aggregator, and receive the updated global model.}
\label{fig:fedavg_round}
\end{figure}

The aggregation step can be expressed as:
\begin{equation}
w_{t+1} = \sum_{k=1}^{K} \frac{n_k}{n} w_{t+1}^{k}
\end{equation}

where $K$ is the number of clients, $n_k$ the number of samples at client $k$, and $n$ the total number of samples.

\subsection{Prototype Configuration}

The current prototype simulates $K=3$ edge nodes, each holding a disjoint partition of the cleaned MQTT dataset. Local models are \texttt{SGDRegressor} instances trained for 5 local epochs per round. Aggregation runs for 10 communication rounds. Although this is a simplified setting, it demonstrates that the global model's error decreases monotonically across rounds, validating the FedAvg workflow without requiring raw data to leave the simulated edge.
